\documentclass[10pt,a4paper]{amsart}
\usepackage[margin=19mm]{geometry}
\usepackage{amsmath,amssymb,mathtools}
\usepackage[scr=rsfs,cal=euler]{mathalfa}
\usepackage{xfrac}
\usepackage[unicode,hidelinks,bookmarks=false]{hyperref}
\newcommand{\ie}{i.e.\ }
\newcommand{\cf}{cf.\ }
\newcommand{\resp}{respectively\ }
\newcommand{\Subsection}{\subsection}
\providecommand{\nobreakdash}{\nobreak\hbox{-}}
\setlength{\emergencystretch}{2em}
\multlinegap0pt
\usepackage[shortlabels]{enumitem}
\usepackage{amssymb}
\usepackage{mathtools}
\usepackage{dsfont}
\usepackage[all]{xy}
\newtheorem{lemma}{Lemma}
\newtheorem{corollary}{Corollary}
\newtheorem{proposition}{Proposition}
\def\I{{\mathcal I}}
\def\M{{\mathcal M}}
\DeclareMathOperator{\depth}{depth}
\DeclareMathOperator{\h}{H}
\newcommand{\m}{\mathfrak m}
\title{Multiplicity versus Buchsbaumness of the special fiber ring}
\author{Anoot Kumar Yadav, Kumari Saloni}
\date{Source: arXiv:2102.04218v4 (CC BY 4.0)}
\begin{document}
\maketitle

\noindent\emph{Source and licence.} Multiplicity versus Buchsbaumness of the special fiber ring. Version arXiv:2102.04218v4 (CC BY 4.0), distributed by arXiv under the Creative Commons Attribution 4.0 International licence, http://creativecommons.org/licenses/by/4.0/, as recorded in the arXiv licence metadata for this exact version and shown on the abstract page https://arxiv.org/abs/2102.04218v4. Full licence notice: \texttt{licenses/arxiv-2102.04218-CC-BY-4.0.txt}.\par\smallskip

\noindent\textit{Editorial scope.} Counted unit: the article's proposition theorem-4.3 (source0.tex lines 901--904). The article's complete original proof of it is delivered with it (source0.tex lines 905--935, one \texttt{proof} environment of 31 lines). No mathematics is written, repaired or re-derived: the statement and the proof are the article's own lines, in the article's own notation, and no retained line is substituted or re-flowed.  Retained uncounted as the source-local context the proof needs: lines 241--242; lines 576--586; lines 847--854.  Nothing was omitted from any retained span, so the package carries no omission record.  No retained line contains a citation, so the wrapper carries no bibliography and no bibliography entry.  No retained line contains a figure environment, an image-inclusion command or drawing code, so no image is delivered and the package declares no asset dependency.  Editorial formatting only, in addition: an AMS \texttt{amsart} preamble that restates the article's own declarations and macros used by the retained lines, namely the environments \texttt{lemma}, \texttt{corollary}, \texttt{proposition} on separate counters, together with the standard packages those lines need.  The retained environments print in this document as Lemma 1, Corollary 1, Proposition 1 in order of appearance, whereas the article prints the same items with the numbering of its own full document.  The complete native source of the article as distributed in the arXiv e-print for this exact version is bundled unchanged as \texttt{sources/107880/source0.tex}.
	 \begin{equation}\label{very-1st-eqn(a)}
	 	e_1(\M)-e_1^J(M)\geq 2e_0(\M)-2\ell(M/M_1)-\ell(M_1/(M_2+JM)) \end{equation}

\begin{lemma}\label{lemma-reduction-steps}
	\begin{enumerate}[label=(\roman*)]
		\item  If $(C_1)$ and $(C_2)$ conditions are satisfied by $a_1,\ldots,a_d$ for $M$, then $(C_1)$ and $(C_2)$ conditions are satisfied by $a_1,\ldots,a_d$ for $\overline{M}$. 
		
		\item \label{lemma-4.1} Suppose that $d\geq 2$. Let $a\in J\setminus \m J$ be a superficial element for $\M$ and $\M_J$.  Then  $2e_0(\M)-e_1(\M)+e_1^J(M)=2\ell_A(M/M_1)+\ell_A(M_1/(M_2+JM))$ if and only if
		$2e_0(\M^\prime)-e_1(\M^\prime)+e_1^J(M^\prime)=2\ell_{A}(M^\prime/{M}^\prime_1)+\ell_{A}({M}^\prime_1/({M}^\prime_2+JM^\prime)).$
		
		\item \label{lemma-4.2} $2e_0(\M)-e_1(\M)+e_1^J(M)=2\ell_A(M/M_1)+\ell_A(M_1/(M_2+JM))$ if and only if  
		$2e_0(\overline{\M})-e_1(\overline{\M})+e_1^J(\overline{M})=2\ell_A(\overline{M}/\overline{M_1} )+\ell_A(\overline{M_1} /(\overline{M_2} +J\overline{M}))$ and $W\subseteq M_2+JM$.
	\end{enumerate}
\end{lemma}

\begin{corollary}\label{corollary-3.6} Let $M$ be a finitely generated faithful flat $A$-module and $\M=\M_L$ or $\M=\M_\I.$ Suppose the conditions $(C_1)$ and $(C_2)$  are satisfied by $a_1,\ldots,a_d$ for both $A$ and $M.$ Also, assume that $C_3$ is satisfied by $A$ and $M$ and  $e_1(\M)-e_1^J(M)=2e_0(\M)-2\ell(M/M_1)-\ell(M_1/(M_2+JM))$ holds. Then the following conditions hold:
	\begin{enumerate}[label=(\roman*)]
		\item $M_{n+2}\subseteq J^n M_1$ for all $n\geq 0, $
		\item $J^nM\cap M_{n+1}=J^n M_1$ for all $n\geq 0,$
		\item $\depth G(\mathcal{\M})>0$ and 
		\item $(a_1,\ldots,\check{a_i},\ldots,a_d)M:_Ma_i\subseteq M_2+JM$ for all $1\leq i\leq d.$
	\end{enumerate}
\end{corollary}

\begin{proposition}\label{theorem-4.3}
 Let $M$ be a finitely generated faithful flat $A$-module and $\M=\M_L$ or $\M=\M_\I$. Suppose $(C_1)$ and $(C_2)$ conditions are satisfied by $a_1,\ldots,a_d$ for both $A$ and $M$.  Let the equality
	$e_1(\M)-e_1^J(M)=2e_0(\M)-2\ell_A(M/M_1)-\ell_A(M_1/(M_2+JM))$ hold. Then $M_{n+2}\subseteq J^nM_2+W$ for $n\geq 1$ and $W\subseteq M_2$.
\end{proposition}

\begin{proof}
 We first show that $W\subseteq M_2$.  By Lemma \ref{lemma-reduction-steps}(iii), we have that $W\subseteq M_2+JM$ and the equality \eqref{very-1st-eqn(a)} holds in $\overline{M}$. Then by Corollary \ref{corollary-3.6}, $J^n\overline{M}\cap \overline{M_{n+1}}=J^n\overline{M_1}$ for $n\geq 0$. 
 Now let $w=i+q\in W$ where $i\in M_2$ and $q\in JM$. Then $q=w-i\in (JM+W)\cap (M_2+W)=JM_1+W$.
	So, $q\in (JM_1+W)\cap JM=JM_1$ since $JM\cap W=0$ which implies $w=i+q\in M_2$. Hence $W\subseteq M_2$.
	
	For the rest of the proof, we may pass to the ring $\overline{A}$, module $\overline{M}$ and assume that the conditions $(C_1)$ and $(C_2)$ are satisfied by the images of $a_1,\ldots,a_d$ in $\overline{A}$ for both $\overline{A}$ and $\overline{M}.$ Moreover,  $\overline{A}$ and $\overline{M}$ both satisfy $(C_3)$. 
	We show, by induction on $d$, that $M_{n+2}= J^n M_2$ for all $n\geq 1$. Suppose the equality  $e_1(\M)-e_1^J(M)=2e_0(\M)-2\ell_A(M/M_1)-\ell_A(M_1/(M_2+JM))$
	holds. Then
	by Corollary \ref{corollary-3.6}, $\depth G(\M)>0$ and $M_{n+2}\subseteq J^n M_1$ for all $n\geq 0.$
	% % %  {\color{blue}{ We need $Q+W\varsubsetneq I_1+W $ in the statement of the theorem.}} 
	We may assume that 
	$f_1=a_1t\in R(J)$ is a regular element of $G(\M)$.
	
	Suppose $d=1$. Since $M_{n+2}\subseteq J^n M_1\subseteq a_1^nt^nM$ and $f_1^n$ is a regular element of $G(\M)$,  we get that $M_{n+2}=M_{n+2}\cap a_1^nt^nM=a_1^nt^nM_2= J^nM_2$ for all $n\geq 1$.
	
	Suppose $d\geq 2.$ As earlier, let $M^\prime=M/a_1M,$ 
	$\M^\prime=\{M_n^\prime=(M_n+a_1M)/a_1M\}_{n\geq0}$ and $\M_J^\prime=\{(J^nM+a_1M)/a_1M\}_{n\geq 0}$ be the filtration of $M^\prime$. Then the conditions  $(C_1)$ and $(C_2)$ are satisfied for $M^\prime$ and
	$e_1(\M^\prime)-e_1^J(M^\prime)=2e_0(\M^\prime)-2\ell_A(M^\prime/{M_1}^\prime)-\ell_A({M_1}^\prime/({M_2}^\prime+JM^\prime))$ by Lemma \ref{lemma-reduction-steps}.
	Further, $(C_1)$, $(C_2)$ and $(C_3)$ are satisfied for $\overline{M^\prime}=M^\prime/\h_{\m}^0(M^\prime)$. 
	Let $\overline{\M^\prime} =\{(M_n^\prime+\h_{\m}^0(M^\prime))/\h_{\m}^0(M^\prime)\}$ be the $I^\prime \overline{A^\prime}$-good filtration of $\overline{M^\prime}$, where $\overline{A^\prime}=A^\prime/\h_{\m}^0(A^\prime)$. 
	By induction hypothesis, we have for all $n\geq 1$,
	\begin{eqnarray*}
		&M_{n+2}^\prime&\subseteq  J^n M_2 ^\prime+\h_{\m}^0(M^\prime)\\
		\Rightarrow & M_{n+2} & \subseteq J^n M_2+a_1M+\h_{\m}^0(M^\prime)\\
		&& \subseteq J^n M_2+(a_1M:J)
	\end{eqnarray*}
	since $M_{n+2}\subseteq J^nM_1$ by Corollary \ref{corollary-3.6}, $a_1,\ldots,a_d$ is a d-sequence of $M$ and $\h_{\m}^0(M^\prime)=(a_1M:J)$. Now we show that $M_{n+2}\subseteq J^nM_2$ by induction on $n$. let $x\in M_{n+2}$, write $x=y+z$ where $y\in J^nM_2$ and 
	$z\in (a_1M:J)$. Then $z=x-y\in (a_1M:J)\cap M_{n+2}\subseteq (a_1M:J)\cap J^n M_1\subseteq (a_1M:J)\cap JM =a_1M$ since
	$M_{n+2}\subseteq J^nM_1$ by Corollary \ref{corollary-3.6} and $a_1,\ldots,a_d$ is a d-sequence of $M$.  This implies $$z\in a_1M\cap M_{n+2}=a_1M_{n+1}\subseteq JM_{n+1}$$ as $a_1t$ is a regular element for $G(\mathcal{M})$. 
	For $n=1$, we get that $x=y+z\in JM_2$. Suppose $n>1$. Then $z\in JM_{n+1}\subseteq J\cdot J^{n-1}M_2$ by induction hypothesis which implies $x\in J^nM_2$. This completes the proof.
\end{proof}

\end{document}
